%---------- Inleiding ---------------------------------------------------------

\section{Inleiding}%
\label{sec:inleiding}

Binnen moderne softwareontwikkeling vormen end-to-end (E2E) testen een onmisbaar pijler voor kwaliteitsborging. Deze testen zijn de belangrijkste methode om te garanderen dat complexe systemen blijven functioneren zoals de eindgebruiker dat verwacht. E2E-testen verifiëren namelijk de volledige gebruikersstroom van begin tot einde, waardoor het inzichten biedt in de productiekwaliteit en de schaal van problemen.
\bigskip
\\
Traditioneel vereisen E2E testen opzetten en onderhouden een continue inbreng van ontwikkelaars, hierdoor is hun kostbare tijd weggenomen van nieuwe functionaliteiten te ontwikkelen. Dus is het belangrijk om een systeem te ontwerpen dat op een zelfvoorzienende manier de productkwaliteit meet bij elke oplevering naar de GitHub repository met minimale belasting voor het ontwikkelingsteam. 
\bigskip
\\
De kwaliteit van het product dat talentguide oplevert is nu laag, zelfs in de productieomgeving voor de eindgebruiker. Om de impact hiervan te begrijpen, moet eerst duidelijk zijn wat het platform precies doet. Talentguide is een SaaS-platform dat organisaties ondersteunt bij het structureren en interpreteren van personeelsgegevens. Het platform verwerkt data zoals cv’s, functiebeschrijvingen en interne documenten en zet deze om naar een gestructureerd competentieprofiel. Daarbij worden eerst kleinere bouwstenen geëxtraheerd, zoals taken, tools, kennis en gedragskenmerken. Binnen talentguide worden deze aangeduid als TTK’s. Deze TTK’s krijgen een expertiselevel en worden vervolgens gegroepeerd onder bredere vaardigheden. Op die manier ontstaat er een profiel dat niet enkel weergeeft welke vaardigheden een werknemer bezit, maar ook op welk niveau deze vaardigheden aanwezig zijn.
\bigskip
\\
Deze profielen vormen de basis voor verdere inzichten binnen het platform. Talentguide gebruikt de gestructureerde data om werknemers te matchen met rollen, competentie gaten te identificeren en ontwikkelingsnoden zichtbaar te maken. Het platform kan ook gelijkaardige of ontbrekende competenties voorstellen om een profiel verder te vervolledigen. Omdat deze inzichten afhankelijk zijn van correcte en betrouwbare data, is de kwaliteit van de onderliggende software belangrijk. Wanneer releases fouten bevatten of gebruikersstromen niet correct werken, kan dit rechtstreeks invloed hebben op de betrouwbaarheid van competentieprofielen, validatieflows en de gebruikerservaring. Daarom is het relevant om te onderzoeken hoe een geautomatiseerde end-to-end testingworkflow kan bijdragen aan een stabieler en betrouwbaarder SaaS-platform.
\bigskip
\\
Het grote probleem is dat er nu geen metingen zijn uitgevoerd op de productkwaliteit, waardoor er geen maatstaf is om te zien hoe groot of klein problemen zijn. Dit komt door het gebrek aan E2E testen in de huidige React/Node.js codebase. Dit leidt ertoe dat releases naar de productieomgeving breken en een slechte gebruikerservaring opleveren.
\bigskip
\\
Om deze productkwaliteit te verankeren zonder de snelheid van het ontwikkelingsteam te belemmeren is de integratie van Artificiële Intelligentie (AI) onderzocht. Hierbij staat de inzet van een geoptimaliseerd vaardigheidsbestand centraal. Door AI-agents via dit bestand te voorzien van technische richtlijnen over de codebase en de structuur van volwaardige testen, kan de drempel voor het genereren van testen worden verlaagd. Deze geautomatiseerde testen zijn vervolgens verankerd in het ontwikkelingsproces door die te integreren binnen de CI/CD pijplijn van de GitHub omgeving. Hierdoor wordt bij elke oplevering van nieuwe functionaliteiten direct gevalideerd, wat een continue terugkoppeling tussen ontwikkeling en kwaliteit garandeert.
\bigskip
\\
Ondanks de efficiëntievoordelen brengt het gebruik van Large Language Models (LLMs) in Quality Assurance (QA) processen specifieke risico's met zich mee, zoals gehallucineerd gebruikersgedrag waarbij de AI foutieve scenario’s genereert om het successtadium te bereiken. Bovendien kan de AI gegenereerde code het debuggen bemoeilijken, waarbij de analyse van testfouten transformeert van het lezen van code naar ingewikkeld "detectivewerk".
\bigskip
\\
De doelgroep van dit onderzoek is talentguide in het bijzonder het development- en QA team. Die hebben behoefte aan een systeem dat hun inzichten geeft op de schaal van de kwaliteitsproblemen, en dat de drempel van E2E testen verlaagd. Dit onderzoek richt zich vooral op het ontwerpen van een E2E testsysteem dat geïntegreerd is binnen de bestaande GitHub- en CI/CD infrastructuur van talentguide, om met minimale intrusie van ontwikkelaars de huidige en toekomstige functionaliteiten continu te kunnen monitoren. 
\bigskip
\\
De hoofdonderzoeksvraag luidt: hoe kan een geautomatiseerd end-to-end testsysteem worden opgezet dat het product continu meet, zodat kwaliteit is verankerd in het ontwikkelingsproces en ervoor zorgt dat ontwikkelaars geen kostbare tijd moeten nemen om volledig nieuwe testen te schrijven?
\bigskip
\\
Om het probleem beter te begrijpen, zijn de volgende aspecten in kaart gebracht:
\begin{itemize}
    \item Wat is de gemeten foutmarge?
    \item Waar bevindt zicht de laagste kwaliteit?
    \item Hoe komt het dat de productkwaliteit laag is?
    \item Hoe kan kwaliteitscontrole een deel worden van het ontwikkelingsproces?
\end{itemize}
\bigskip
Om een oplossing te vinden, zijn de volgende aspecten onderzocht:
\begin{itemize}
    \item Hoe kunnen geautomatiseerde E2E testen opgezet worden?
    \item Welke procedures kunnen opgezet worden om de productkwaliteit te verhogen?
    \item Kan er gebruik worden gemaakt van AI?
    \item Welke tools bestaan er om E2E testen te automatiseren?
\end{itemize}
\bigskip
Het doel van deze bachelorproef is om de kwaliteit van het eindproduct bij talentguide te verhogen door een functioneel prototype af te leveren van een geautomatiseerde testingworkflow dat op basis van minimale ontwikkelaars input robuuste testen genereert en uitvoert.
\bigskip
\\
Het eindresultaat wordt als een succes beschouwd wanneer dit systeem objectief de productkwaliteit in kaart brengt bij het opleveren van nieuwe functionaliteiten naar de Github repository. Hiermee is de fundamentele stap gezet om structurele kwaliteitsmetingen binnen de organisatie te introduceren.


%---------- Stand van zaken ---------------------------------------------------

\section{Literatuurstudie}%
\label{sec:literatuurstudie}

\section{E2E-testen in de ontwikkelingscyclus}

\subsection{Waarom E2E-testen nodig zijn}

Voor webapplicaties is de browserinterface een belangrijke plaats waar fouten zichtbaar worden voor eindgebruikers \autocite{DiMeglio2026}. Daarom volstaat het niet altijd om enkel afzonderlijke functies, componenten of API's te testen. E2E-testen controleren of volledige gebruikersstromen werken zoals verwacht, bijvoorbeeld door een applicatie via de browser te bedienen en daarbij navigatie, invoer, interacties en validaties te doorlopen. Web-GUI-testen sluiten aan bij die benadering, omdat deze testvorm de kwaliteit van een webapplicatie beoordeelt vanuit het perspectief van de eindgebruiker \autocite{DiMeglio2026}.

\textcite{DiMeglio2026} tonen aan dat web-GUI-testen in de praktijk worden toegepast met browserautomatiseringsframeworks zoals Selenium, Playwright, Cypress en Puppeteer. Deze frameworks maken het mogelijk om gebruikersinteracties herhaalbaar uit te voeren en te controleren of een applicatie zich correct gedraagt in realistische scenario's. Voor een SaaS-platform zoals talentguide sluit dat rechtstreeks aan bij productkwaliteit: afzonderlijke technische onderdelen moeten werken, maar gebruikers moeten ook volledige workflows betrouwbaar kunnen doorlopen.

\subsection{E2E-testen als bewaking van gebruikersstromen}

Onderzoek naar web-GUI-testen toont het belang aan van testen die complex UI-gedrag op applicatieniveau controleren \autocite{DiMeglio2026}. De waarde van E2E-testen ligt vooral in het bewaken van samenhang tussen verschillende onderdelen van een applicatie. Waar unit- en integratietesten specifieke functies of technische koppelingen controleren, richten E2E-testen zich op het gedrag van de applicatie als geheel. Dat maakt deze testvorm geschikt om gebruikersstromen te bewaken die meerdere schermen, componenten, invoervelden en achterliggende services combineren.

\textcite{DiMeglio2026} tonen daarnaast dat web-GUI-testen in open-sourceprojecten kunnen meegroeien met de applicatiecode. Hun analyse van 472 webapplicaties toont een verband met meer ontwikkelactiviteit, meer contributors en actievere issue-trackers. Een rechtstreeks kwaliteitsvoordeel volgt daar niet uit. Wel toont de studie dat zulke testen passen binnen actieve ontwikkeling. In dit onderzoek worden E2E-testen daarom gezien als een manier om productkwaliteit al tijdens de ontwikkeling zichtbaar te maken.

\subsection{Waarom E2E-testen onderhoud vragen}

\textcite{DiMeglio2026} beschrijven dat scripted web-GUI-testen technische kennis vragen, omdat testcode geschreven en aangepast moet worden wanneer de applicatie verandert. De studie onderscheidt scripted testing van low-code-, capture-and-replay- en scriptless-benaderingen. Scripted testing biedt meer controle over selectors, assertions, testdata en teststructuur, maar vraagt ook meer programmeerkennis en onderhoud.

Een belangrijk onderhoudsrisico bij web-GUI-testen is fragility: testen kunnen breken door kleine wijzigingen in de interface, de DOM-structuur of gebruikte selectors \autocite{DiMeglio2026}. Een ander risico is flakiness: testen falen dan niet-deterministisch, bijvoorbeeld door timingproblemen, asynchrone inhoud of wisselende testdata \autocite{DiMeglio2026}. Zulke problemen zijn belangrijk binnen een ontwikkelingsproces, omdat een foutief falende test het werk vertraagt, terwijl een foutief geslaagde test een vals gevoel van kwaliteit kan geven.

De studie van \textcite{DiMeglio2026} nuanceert tegelijk het idee dat web-GUI-testen per definitie onstabiel zijn. Hun resultaten tonen dat web-GUI-testen in veel projecten relatief lang kunnen blijven bestaan en vaak slechts beperkt aangepast worden. Dat wijst erop dat E2E-testen onderhoudbaar kunnen zijn wanneer het ontwerp goed is en de testen stapsgewijs worden bijgewerkt. De uitdaging ligt dus niet enkel bij het schrijven van E2E-testen. Ook de testaanpak moet betrouwbaar, onderhoudbaar en toepasbaar blijven binnen de normale ontwikkelingscyclus.



% Voor literatuurverwijzingen zijn er twee belangrijke commando's:
% \autocite{KEY} => (Auteur, jaartal) Gebruik dit als de naam van de auteur
%   geen onderdeel is van de zin.
% \textcite{KEY} => Auteur (jaartal)  Gebruik dit als de auteursnaam wel een
%   functie heeft in de zin (bv. ``Uit onderzoek door Doll & Hill (1954) bleek
%   ...'')


%---------- Methodologie ------------------------------------------------------
\section{Methodologie}%
\label{sec:methodologie}

\subsection{Fase 1:  Nulmeting van de huidige staat}
Deze fase geeft verduidelijking op de volgende deelvragen:
\begin{itemize}
    \item Wat is de gemeten foutmarge?
    \item Waar bevindt zicht de laagste kwaliteit?
    \item Hoe komt het dat de productkwaliteit laag is?
\end{itemize}

De eerste fase richt zich op het in kaart brengen van de oorzaken achter de huidige kwaliteitsgebreken binnen het SaaS-platform. In eerste instantie zal worden onderzocht waarom de bestaande kwaliteitscontroles falen tijdens het integratieproces; hierbij wordt de GitHub CI-pijplijn geanalyseerd om de technische blokkades te identificeren die voorkomen dat testen succesvol worden uitgevoerd bij een push naar de repository.
\bigskip
\\
Vervolgens zal het Sentry dashboard kritisch worden onderzocht, omdat dit platform elke foutmelding in de development, staging en productieomgeving monitort en logt. Door deze incidenten te onderzoeken kan een beeld gevormd worden van de manier waarop deze defecten de gebruikerservaring van de eindgebruiker aantasten. De vijf meest significante fouten worden hierna diepgaand onderzocht om de exacte bron van het conflict te achterhalen.
\bigskip
\\
Na deze fase is er inzicht in de tekortkomingen in de huidige opzet, wat als basis dient voor het ontwerpen van een oplossing die is toegespitst op de geïdentificeerde noden.

\subsection{Fase 2: Richtlijnen en benodigheden}
Deze fase geeft verduidelijking op de volgende deelvragen:
\begin{itemize}
    \item Hoe kan kwaliteitscontrole een deel worden van het ontwikkelingsproces?
    \item Kan er gebruik worden gemaakt van AI?
    \item Welke tools bestaan er om E2E testen te automatiseren?
\end{itemize}

Nadat de oorzaken voor de huidige kwaliteit tekortkomingen zijn geïdentificeerd, richt deze fase zich op de richtlijnen en benodigdheden te definiëren. van de AI testworkflow. Het doel hiervan is het opstellen van de nodige ricthlijnen voor elk vaardigheidsbestand, hoe de keten samenwerkt, wat de test coverage niveaus zijn en de functionele en non-functionele technische richtlijnen voor elk vaardigheidsbestand. Dit waarborgt dat de AI output artefacten oplevert die kunnen integreren in de codebase van talentguide.
Deze eisen garanderen de schaalbaarheid en toekomstbestendigheid van de workflow. In overleg met de IT lead en de copromotor is besloten om Playwright te gebruiken als het testing raamwerk voor de oplossing.
\bigskip
\\
Na deze fase zijn er tabellen met vereisten en kwaliteitsnormen die als basis dienen voor de ontwikkeling van het prototype.


\subsection{Fase 3: Opzet van het prototype}
Deze fase geeft verduidelijking op de volgende deelvragen:
\begin{itemize}
    \item Hoe kunnen geautomatiseerde E2E testen opgezet worden?
    \item Welke procedures kunnen opgezet worden om de productkwaliteit te verhogen
\end{itemize}

In deze fase zal er een proof-of-concept woorden ontwikkeld volgens de richtlijnen die zijn vastgesteld in de vorige fase. Dit prototype bestaat uit een geautomatiseerde testingworkflow die integreert met de ontwikkelingscyclus van talentguide. De opzet start met het invoeren van het Playwright framework, waarbij de configuratie zal worden afgestemd op de eisen van de repository en CI/CD pijplijn.
\bigskip
\\
Bij het schrijven van de testen worden de beste praktijken van Playwright gebruikt. Hier zal worden gebruikgemaakt van het Page Object Model (POM). Hierbij zullen de interacties met de gebruikersinterface worden geabstraheerd in afzonderlijke klassen. Daarnaast zal er worden gewerkt met fixtures en een centrale support folder voor mockdata, constanten en helper methodes. Voor de authenticatie zal de officiële Playwright auth-setup geïmplementeerd worden, hierdoor kan de inlog sessie gedeeld worden over de volledige testsuite.
\bigskip
\\
Ook is een noodzakelijk onderdeel van deze automatisering het ontwikkelen van robuuste vaardigheidsbestanden. Dit document dient als technische richtlijnen voor de AI agents om met minimale tussenkomst van de ontwikkelaar end-to-end testen te kunnen generen. Door deze toevoeging zal de drempel van Test Driven Development (TDD) verlaagd worden en wordt de objectieve meting van d productkwaliteit gefaciliteerd.
\bigskip
\\
Na deze fase is er een praktische blauwdruk voor toekomstige uitvoering met een coding agent.


\subsection{Fase 4: Evaluatie}

Deze fase richt zich op de beoordeling van het prototype door het evaluatie vaardigheidsbestand in te zetten als een objectieve maatstaf voor de kwaliteitsborging. Hierbij is het kwaliteitsniveau toegekend a.d.h.v minimale criteria en de toepasbaarheid van de technische richtlijnen.
\bigskip
\\
Na deze fase kan een conclusie gevormd worden als de ontwikkelde aanpak de productkwaliteit daadwerkelijk kan verbeteren.

%---------- Verwachte resultaten ----------------------------------------------
\section{Verwacht resultaat, conclusie}%
\label{sec:verwachte_resultaten}

Het verwachte resultaat van dit onderzoek is het identificeren waar de laagste productkwaliteit zich bevindt, om vervolgens een structurele oplossing in te voeren die de foutmarge binnen het platform verlaagt. Dit proces vertrekt vanuit een grondige analyse van de vijf meest voorkomende incidenten in Sentry. Het opgeleverde artefact is een werkend prototype van een geautomatiseerde end-to-end testingworkflow aangedreven door Playwright en ondersteund door een AI-gestuurde workflow. De kern van deze automatisering bestaat uit een vaardigheidsbestand die de nodige technische richtlijnen bevat voor de generatie van testen.
De centrale hypothese is dat de integratie van deze vaardigheidsbestanden de drempel voor Test Driven Development (TDD) verlaagt, waardoor AI-agents in staat zijn om met minimale menselijke tussenkomst kwalitatieve testen te genereren. Daarnaast zal het prototype meetbare data opleveren over de feedbackloop van de CI/CD pijplijn en de kostenefficiëntie van de gekozen methode.

\subsection{Conclusie}

De meerwaarde van dit onderzoek is om een robuuste testprocedure te verankeren in het ontwikkelingsproces. Hierdoor zal de productieomgeving niet dienen als testgrond van het systeem door de eindgebruiker. Voor het development team betekent dat er een tool aanwezig is die de drempel verlaagd voor de kwaliteitsborging. Voor talentguide betekent dit dat de betrouwbaarheid van het SaaS-platform is verhoogd, wat ervoor kan zorgen dat de klantentevredenheid toeneemt.


